% !TeX program = xelatex
\documentclass[UTF8,11pt,fleqn,fontset=none]{ctexart}

\setCJKmainfont{Songti SC}
\setCJKsansfont{PingFang SC}
\setCJKmonofont{PingFang SC}

\usepackage[margin=1in]{geometry}
\usepackage{amsmath}
\usepackage{amssymb}
\usepackage{amsthm}
\usepackage{mathtools}
\usepackage{bm}
\usepackage{graphicx}
\usepackage{hyperref}

\renewcommand{\Im}{\mathrm{Im}}
\renewcommand{\Re}{\mathrm{Re}}
\newcommand{\ii}{\mathrm{i}}
\newcommand{\ee}{\mathrm{e}}
\newcommand{\dd}{\mathrm{d}}
\newcommand{\CC}{\mathbb{C}}
\newcommand{\RR}{\mathbb{R}}

\theoremstyle{definition}
\newtheorem{conjecture}{猜想}[section]
\newtheorem{assumption}[conjecture]{假设}
\newtheorem{proposition}[conjecture]{命题}
\newtheorem{theorem}[conjecture]{定理}

\title{周期线缺陷波导特征值的边界积分计算与后验误差估计}
\author{}
\date{\today}

\begin{document}

\maketitle

\section{引言}
\label{sec:introduction}

二维光子晶体中的线缺陷破坏了横向的完整周期结构，却保留了沿缺陷方向的周期性。
当频率落在背景周期介质的谱隙内时，这种缺陷可能支持沿缺陷传播而在横向局域的导模。
相应问题既不是完整周期介质的普通能带计算，也不是均匀开放波导的横截面特征值问题；
它是在无界周期背景中寻找缺陷诱导的离散谱。Fliss 利用准周期谱分解和
Dirichlet-to-Neumann 映射，把这类导模刻画为固定 Bloch 波数下的条带特征值问题
\cite{Fliss2013}。

开放周期端的精确处理和高精度离散是导模计算的两个核心困难。边界积分方法能够只离散
材料界面，并已被用于构造高阶波导模态求解器\cite{LuLu2014}；但一个很小的离散矩阵残量
并不自动给出连续物理特征值的可靠位数。Barnett、Hassell 与 Tacy 以边界量控制谱距离的
工作说明，边界上的不相容量可以成为谱包含的有效入口\cite{BarnettHassellTacy2018}；不过
其有界光滑 Neumann 问题与本文的无界周期传输问题不同，不能直接给出本文的结论。光子晶体
特征值问题中的有限元研究也表明，后验估计可以用于衡量特征值误差并指导离散加密
\cite{EngstromGianiGrubisic2016}。本文为当前条带算子另外证明一个预解式谱距离定理，并用
经过向外舍入的区间算术给出各项输入的可靠单侧界。

本文研究两侧为周期介质、中心为一般局部缺陷的二维波导。每个含材料的基本胞元至多含一个
单连通介质柱截面；中心胞元可以含一个形状不同的介质柱，也可以为空。后者只是本文数值
算例采用的特例，而不是一般模型的限制。与把材料写成连续折射率函数的通常表述相比，本文
采用具有尖锐界面的分片常数折射率，因此偏微分方程必须在各均匀子区域内理解，并在每条
材料界面上显式满足传输条件。

本文首先在固定 Bloch 波数下用边界积分、端口 Fourier 截断和有序 QZ 分解计算候选
$\widehat k$ 及其分片重构场。随后固定同一候选、密度、传播矩阵、材料和几何输入，直接
评价原始场在全部材料界面与人工拼接界面上的值跳跃和普通法向导数跳跃。正移位预解式的
全界面迹估计把这些跳跃转化为
$\operatorname{dist}(\widehat k^2,\sigma(A(\beta)))$ 的可计算上界；这一全谱距离结论本身
不要求事先验证本质谱隙。只有另行证明所得包含区间落在同一算子的本质谱隙内，才能进一步
把其中的谱点解释为离散导模特征值。本文不据此声称特征值唯一、编号已知或谱隙本身已经
得到证明。第~\ref{sec:modelproblem}~节给出模型，第~\ref{sec:numerics}~节给出候选算法，
第~\ref{sec:error}~节证明谱距离界，第~\ref{sec:results}~节报告一个圆形介质柱、空中心的
谱距离可靠上界的计算实例。


\section{模型问题}
\label{sec:modelproblem}

首先在整个平面 $\RR^2$ 上定义物理问题。设 $\rho_{\mathrm p}$ 描述未引入缺陷时的
二维光子晶体。本文假定 $\rho_{\mathrm p}$ 为正的分片常数函数，并且分别以
$L>0$ 和 $d>0$ 为 $x$、$y$ 两个坐标方向的周期，即存在常数
$0<\rho_-\leq\rho_+<\infty$，使得
\begin{align}
  \rho_-&\leq \rho_{\mathrm p}(x,y)\leq \rho_+,
  &
  \rho_{\mathrm p}(x+jL,y+md)
  &=\rho_{\mathrm p}(x,y)
  \label{eq:periodic-background}
\end{align}
对几乎处处的 $(x,y)\in\RR^2$ 以及所有 $j,m\in\mathbb Z$ 成立。沿 $y$ 方向引入一条
无限长线缺陷。按照左侧背景、中心缺陷和右侧背景的宏观分区，记
\begin{align}
  \mathcal{W}_-&=(-\infty,-a)\times\RR,
  &\mathcal{W}_0&=(-a,a)\times\RR,
  &\mathcal{W}_+&=(a,\infty)\times\RR,
  \qquad a>0.
  \label{eq:defect-region}
\end{align}
其中 $\mathcal{W}_0$ 是整个线缺陷区域。缺陷区中的材料系数 $\rho_0$ 同样取为正的
分片常数函数，并保留 $y$ 方向的周期性：
\begin{align}
  \rho_-&\leq \rho_0(x,y)\leq\rho_+,
  &
  \rho_0(x,y+md)&=\rho_0(x,y),
  \qquad (x,y)\in\mathcal{W}_0,\quad m\in\mathbb Z.
  \label{eq:periodic-defect}
\end{align}
若背景与缺陷原先具有可公度的纵向周期，则将 $d$ 取为它们的一个公共周期。整个二维
缺陷晶体由
\begin{align}
  \rho(x,y)&=
  \begin{cases}
    \rho_{\mathrm p}(x,y), & (x,y)\in\RR^2\setminus\mathcal{W}_0,\\
    \rho_0(x,y), & (x,y)\in\mathcal{W}_0
  \end{cases}
  \label{eq:full-plane-coefficient}
\end{align}
描述。这个构造与 Fliss 对二维周期介质中线缺陷的几何描述相同
\cite[第~2~节]{Fliss2013}；区别在于本文把材料明确建模为分片常数，因此必须保留尖锐界面
及其传输条件。

记 $\Upsilon$ 为各介质柱截面边界以及这些边界按相应背景或缺陷周期平移后产生的全部材料
界面的并集；若一个胞元含有多根介质柱，则其每个截面边界都计入 $\Upsilon$。给定波数
$k>0$，全平面上的标量场 $w$ 在每个均匀材料子区域内满足
\begin{align}
  -\Delta w&=k^2\rho w,
  \qquad \text{在 }\RR^2\setminus\Upsilon.
  \label{eq:full-plane-wave-equation}
\end{align}
在 $\Upsilon$ 的任一界面上固定同一个单位法向量 $\nu$，并以 $w^+$、$w^-$ 表示界面两侧的
迹，则要求
\begin{align}
  w^+&=w^-,
  &
  \partial_\nu w^+&=\partial_\nu w^-.
  \label{eq:transmission-conditions}
\end{align}
第一条等式表示场连续，第二条等式表示按同一法向取值的法向导数连续。

令
$B=\RR\times(-d/2,d/2)$，并取
$\beta\in[-\pi/d,\pi/d)$。按照 Fliss 的式~(2.2)，波数 $k>0$ 处的一个\emph{导模}
是全平面问题~\eqref{eq:full-plane-wave-equation}--\eqref{eq:transmission-conditions} 的非零解
$w$，它可以写成
\begin{align}
  w(x,y)&=v(x,y)\ee^{\ii\beta y}.
  \label{eq:guided-mode-definition}
\end{align}
这里 $v$ 沿 $y$ 方向以 $d$ 为周期，并且 $v|_B\in H^1(B)$
\cite[(2.2)]{Fliss2013}。条件 $v|_B\in H^1(B)$ 已经包含无界 $x$ 方向上的平方可积性，
因而表达了导模的横向局域性。

记条带的上、下边界分别为
$\Sigma=\RR\times\{d/2\}$ 和 $\widetilde\Sigma=\RR\times\{-d/2\}$，并定义
\begin{align}
  H^1_\beta(B)
  :=\left\{u\in H^1(B):
  u|_\Sigma=\ee^{\ii\beta d}u|_{\widetilde\Sigma}
  \text{，按迹意义}\right\}.
  \label{eq:quasiperiodic-space}
\end{align}
令 $u=w|_B$。与 Fliss 的式~(2.4) 相同，求导模等价于求
$k>0$、$\beta\in[-\pi/d,\pi/d)$ 以及
$u\in H^1_\beta(B)\setminus\{0\}$，使得
\begin{equation}
  \left\{
  \begin{aligned}
    -\Delta u&=k^2\rho u
    &&\text{在 }B\setminus\Upsilon,\\
    u^+&=u^-,
    &\partial_\nu u^+&=\partial_\nu u^-
    &&\text{在 }\Upsilon\cap B,\\
    \partial_yu|_\Sigma
    &=\ee^{\ii\beta d}\partial_yu|_{\widetilde\Sigma}.
  \end{aligned}
  \right.
  \label{eq:guided-mode-strip-problem}
\end{equation}
其中第一行在各均匀材料子区域内成立，第二行使用与
式~\eqref{eq:transmission-conditions} 相同的法向约定；最后一行给出法向导数的准周期条件，
而函数值的准周期条件已经包含在
$u\in H^1_\beta(B)$ 中。这就是本文要求求解的导模问题
\cite[(2.4)]{Fliss2013}。所谓 $\beta$-表述，是先固定实准周期参数 $\beta$，再求
特征值 $k^2$ 及相应的非零场 $u$。

为把这一问题写成算子特征值问题，取加权 Hilbert 空间
\begin{align}
  \mathcal H&=L^2(B,\rho\,\dd x\,\dd y),
  \label{eq:weighted-hilbert-space}
\end{align}
并记
\begin{align}
  H^1(\Delta,B)
  :=\left\{u\in H^1(B):\Delta u\in L^2(B)\right\}.
  \label{eq:H1-Delta}
\end{align}
依照 Fliss 第~3~节开头的定义，令
\begin{equation}
  \begin{aligned}
    A(\beta)&:D(A(\beta))\subset\mathcal H\longrightarrow\mathcal H,
    &
    A(\beta)u&=-\rho^{-1}\Delta u,\\
    D(A(\beta))
    &=\left\{u\in H^1(\Delta,B):
    \begin{array}{l}
      u|_\Sigma=\ee^{\ii\beta d}u|_{\widetilde\Sigma},\\[-1mm]
      \partial_yu|_\Sigma
      =\ee^{\ii\beta d}\partial_yu|_{\widetilde\Sigma}
    \end{array}\right\}.
  \end{aligned}
  \label{eq:fixed-beta-operator}
\end{equation}
这里的 $\Delta u$ 按分布意义理解。由于 $u\in H^1(B)$ 且 $\rho$ 为分片常数，
$\Delta u\in L^2(B)$ 排除了材料界面上的分布源；对于分片光滑的场，这恰好给出
$u^+=u^-$ 和 $\partial_\nu u^+=\partial_\nu u^-$。因此传输条件已经包含在
$D(A(\beta))$ 中。Fliss 证明 $A(\beta)$ 在 $\mathcal H$ 中非负且自伴
\cite[第~3~节及命题~3.1]{Fliss2013}。

于是固定 $\beta$ 的导模问题准确地写成
\begin{align}
  \text{求 }k^2>0\text{ 及 }u\in D(A(\beta))\setminus\{0\},
  \qquad A(\beta)u=k^2u.
  \label{eq:beta-formulation}
\end{align}

为刻画 $A(\beta)$ 的本质谱，先去掉线缺陷，并以周期背景系数
$\rho_{\mathrm p}$ 定义
\begin{align}
  \mathcal H_{\mathrm p}
  &=L^2(B,\rho_{\mathrm p}\,\dd x\,\dd y),
  &
  A_{\mathrm p}(\beta)u
  &=-\rho_{\mathrm p}^{-1}\Delta u,
  &
  D(A_{\mathrm p}(\beta))&=D(A(\beta)).
  \label{eq:periodic-background-operator}
\end{align}
这里最后一个等号指两个定义域由同一组 $y$ 方向准周期条件给出；由于
$\rho$ 与 $\rho_{\mathrm p}$ 均一致正且有界，相应的加权 $L^2$ 范数与通常的
$L^2$ 范数等价。再令
\begin{align}
  C&=(-L/2,L/2)\times(-d/2,d/2)
  \label{eq:periodic-cell}
\end{align}
为周期背景的一个胞元。对
$\theta\in[-\pi/L,\pi/L)$，在
$L^2(C,\rho_{\mathrm p}\,\dd x\,\dd y)$ 中定义胞元 Bloch 算子
\begin{equation}
  \begin{aligned}
    A_{\mathrm p}(\beta,\theta)u
    &=-\rho_{\mathrm p}^{-1}\Delta u,\\
    D(A_{\mathrm p}(\beta,\theta))
    =\left\{u\in H^1(\Delta,C):{}
    \begin{aligned}
      u|_{x=L/2}
      &=\ee^{\ii\theta L}u|_{x=-L/2},\\
      \partial_xu|_{x=L/2}
      &=\ee^{\ii\theta L}\partial_xu|_{x=-L/2},\\
      u|_{y=d/2}
      &=\ee^{\ii\beta d}u|_{y=-d/2},\\
      \partial_yu|_{y=d/2}
      &=\ee^{\ii\beta d}\partial_yu|_{y=-d/2}
    \end{aligned}\right\}.
  \end{aligned}
  \label{eq:background-bloch-fiber}
\end{equation}
这里 $H^1(\Delta,C)=\{u\in H^1(C):\Delta u\in L^2(C)\}$，上述边界等式均按迹意义理解；分布条件
$\Delta u\in L^2(C)$ 仍然包含周期背景内部各材料界面上的传输条件。

\begin{conjecture}[分片常数线缺陷算子的本质谱]
  \label{conj:essential-spectrum}
  对每个固定的 $\beta\in[-\pi/d,\pi/d)$，算子 $A(\beta)$ 的本质谱满足
  \begin{align}
    \sigma_{\mathrm{ess}}(A(\beta))
    &=\sigma(A_{\mathrm p}(\beta))
      =\bigcup_{\theta\in[-\pi/L,\pi/L)}
        \sigma(A_{\mathrm p}(\beta,\theta))
      \notag\\
    &=[0,+\infty)\setminus
      \bigcup_{n=1}^{N(\beta)}(a_n(\beta),b_n(\beta)),
    \qquad
    0\leq N(\beta)\leq+\infty,
    \label{eq:conjectured-essential-spectrum}
  \end{align}
  其中各 $(a_n(\beta),b_n(\beta))$ 是两两不交的开区间，并且
  $0\leq a_n(\beta)<b_n(\beta)$。
\end{conjecture}

猜想~\ref{conj:essential-spectrum} 模仿 Fliss 的命题~3.1
\cite[命题~3.1]{Fliss2013}，但不把原命题直接作为本文模型的已证结论。本文的
$\rho$ 和 $\rho_{\mathrm p}$ 具有尖锐材料界面，算子定义域以分布 Laplacian 和传输条件实现；
因此仍需针对这一分片常数实现核验缺陷扰动的预解式紧性，以及周期背景算子的
Floquet--Bloch 纤维分解。上述适配尚未在本文中完成，所以本质谱等式在此保留为猜想。

\begin{assumption}[存在选定的投影谱隙]
  \label{ass:projected-gap}
  固定 $\beta\in[-\pi/d,\pi/d)$，假设猜想~\ref{conj:essential-spectrum} 成立且
  $N(\beta)\geq1$。选择其中一个严格位于正半轴的非空谱隙，并记为
  \begin{align}
    G_\beta&=(g_-,g_+),
    &0<g_-<g_+<+\infty,
    &\qquad G_\beta\cap\sigma_{\mathrm{ess}}(A(\beta))=\varnothing.
    \label{eq:selected-projected-gap}
  \end{align}
\end{assumption}

以下采用 Weyl--Fredholm 本质谱约定。
\begin{proposition}[谱隙内的离散谱]
  \label{prop:discrete-spectrum-in-gap}
  在假设~\ref{ass:projected-gap} 下，
  $\sigma(A(\beta))\cap G_\beta$ 只由有限重的孤立特征值组成；这些特征值只能在
  $g_-$ 或 $g_+$ 处聚集，在 $G_\beta$ 内部没有聚集点。
\end{proposition}

命题~\ref{prop:discrete-spectrum-in-gap} 是 Fliss 命题~3.3 在当前条件性框架中的对应表述
\cite[命题~3.3]{Fliss2013}。它并不保证
$\sigma(A(\beta))\cap G_\beta$ 非空，只刻画该交集一旦非空时的谱类型。
除非另有说明，后文始终固定满足假设~\ref{ass:projected-gap} 的 $\beta$ 和
$G_\beta$，所有特征值定位、离散化与后验误差估计都只考虑
\begin{align}
  \text{求 }k^2\in G_\beta\text{ 及 }u\in D(A(\beta))\setminus\{0\},
  \qquad A(\beta)u=k^2u,
  \label{eq:gap-restricted-beta-formulation}
\end{align}
亦即只搜索 $k\in(\sqrt{g_-},\sqrt{g_+})$；谱隙外的谱点不属于本文后续的搜索对象。


\section{中心胞元约化与连续 DtN 算子}
\label{sec:dtn}

以下固定第~\ref{sec:modelproblem}~节中的 $\beta$ 以及选定谱隙
$G_\beta=(g_-,g_+)$。沿 $y$ 方向施加 $\beta$-准周期条件以后，条带
$B=\RR\times(-d/2,d/2)$ 在 $x$ 方向仍然无界。记
$I_d=(-d/2,d/2)$。由式~\eqref{eq:defect-region} 的全平面宏观分区定义
\begin{align*}
  B_\varsigma&=\mathcal{W}_\varsigma\cap B,
  \qquad \varsigma\in\{-,0,+\},\\
  B_-&=(-\infty,-a)\times I_d,
  &B_0&=(-a,a)\times I_d,
  &B_+&=(a,\infty)\times I_d.
\end{align*}
中心缺陷胞元及紧邻它的左右周期胞元分别定义为
\begin{align}
  \mathcal{C}_0&=B_0=(-a,a)\times I_d,
  &\mathcal{C}_-&=(-a-L,-a)\times I_d,
  &\mathcal{C}_+&=(a,a+L)\times I_d.
  \label{eq:center-and-lead-cells}
\end{align}
这里 $\mathcal{C}_0$ 是宽度为 $2a$ 的功能性缺陷胞元，其宽度不要求等于 $L$；
$\mathcal{C}_-$ 和 $\mathcal{C}_+$ 才是 $x$ 方向周期为 $L$ 的一个基本背景胞元，而不是
半无限区域。记 $e_x=(1,0)$ 为 $x$ 方向的单位向量。左右半条带可以由这些周期胞元表示为
\begin{align}
  B_-
  &=\operatorname{int}\!\left(
    \bigcup_{j=0}^{\infty}
    (\overline{\mathcal{C}_-}-jL e_x)\right),
  &
  B_+
  &=\operatorname{int}\!\left(
    \bigcup_{j=0}^{\infty}
    (\overline{\mathcal{C}_+}+jL e_x)\right).
  \label{eq:half-guides}
\end{align}
对 $\varsigma\in\{-,0,+\}$，以
$\Omega_\varsigma\subset\mathcal{C}_\varsigma$ 表示相应胞元中的单个介质柱截面区域。
其中 $\Omega_0$ 是中心缺陷胞元中的介质柱区域，并允许
$\Omega_0=\varnothing$；$\Omega_-$ 和 $\Omega_+$ 分别是左右基本周期胞元中的介质柱区域。
中心介质柱若存在则完全包含在 $\mathcal{C}_0=B_0$ 中；$B_-$ 和 $B_+$ 内的介质柱区域分别由
$\Omega_-$ 的平移 $\Omega_--jLe_x$ 和 $\Omega_+$ 的平移 $\Omega_++jLe_x$，
$j\geq0$，生成。因此 $B_\pm$ 中的系数均为周期背景系数 $\rho_{\mathrm p}$，
$\mathcal{C}_0$ 中则保留第~\ref{sec:modelproblem}~节定义的缺陷系数 $\rho$ 及
$\Upsilon$ 所表示的全部尖锐材料界面。人工界面和法向约定为
\begin{align}
  \Gamma_\pm
  &=\partial\mathcal{C}_0\cap\partial B_\pm,
  &\Gamma_-&=\{-a\}\times I_d,
  &\Gamma_+&=\{a\}\times I_d,
  \notag\\
  \nu_-&=-e_x,
  &\nu_+&=e_x.
  \label{eq:artificial-interfaces}
\end{align}
两个法向都从 $\mathcal{C}_0$ 指向相应半条带 $B_\pm$。特别地，它们是中心区域的外法向，
却与 $B_\pm$ 在 $\Gamma_\pm$ 上的外法向相反；下文的 $\partial_{\nu_\pm}$ 始终采用
式~\eqref{eq:artificial-interfaces} 的定向。

对 $U\in\{\mathcal{C}_0,B_-,B_+\}$，只引入后文所需的准周期能量空间
\begin{align}
  H^1_\beta(U)
  :=\left\{z\in H^1(U):
  z|_{y=d/2}=\ee^{\ii\beta d}z|_{y=-d/2}
  \text{，按迹意义}\right\}.
  \label{eq:local-quasiperiodic-space}
\end{align}
令 $\gamma_\pm$ 表示 $\Gamma_\pm$ 上的 Dirichlet 迹，并以迹像及其商范数定义
\begin{align}
  H^{1/2}_\beta(\Gamma_\pm)
  &:=\gamma_\pm H^1_\beta(\mathcal{C}_0)
    =\gamma_\pm H^1_\beta(B_\pm),
  &
  H^{-1/2}_\beta(\Gamma_\pm)
  &:=H^{1/2}_\beta(\Gamma_\pm)^*.
  \label{eq:interface-trace-spaces}
\end{align}
这样，$y$ 方向的 $\beta$-准周期兼容性已经编码在迹空间中，无须借助临界
$H^{1/2}$ 函数的端点点值。记
$\langle\cdot,\cdot\rangle_{\Gamma_\pm}$ 为
$H^{-1/2}_\beta(\Gamma_\pm)$ 与 $H^{1/2}_\beta(\Gamma_\pm)$ 之间的对偶配对。

在每条材料界面上任意固定一个两侧共用的单位法向 $\nu$。现在固定 $k>0$ 且
$k^2\in G_\beta$。给定任一方向和界面数据
$\varphi\in H^{1/2}_\beta(\Gamma_\pm)$，相应的连续半波导 Dirichlet 问题是求
$u_\pm=u_\pm(\varphi;k)\in H^1_\beta(B_\pm)$，使
\begin{equation}
  \left\{
  \begin{aligned}
    -\Delta u_\pm&=k^2\rho_{\mathrm p}u_\pm
    &&\text{在 }B_\pm\setminus\Upsilon\text{ 的各均匀材料子区域内},\\
    u_\pm^+&=u_\pm^- ,
    &\partial_\nu u_\pm^+&=\partial_\nu u_\pm^-
    &&\text{在 }\Upsilon\cap B_\pm,\\
    u_\pm|_{y=d/2}
    &=\ee^{\ii\beta d}u_\pm|_{y=-d/2},
    &\partial_yu_\pm|_{y=d/2}
    &=\ee^{\ii\beta d}\partial_yu_\pm|_{y=-d/2},\\
    u_\pm|_{\Gamma_\pm}&=\varphi.&&
  \end{aligned}
  \right.
  \label{eq:half-guide-dirichlet-problem}
\end{equation}
第二行正是当前分片常数模型的传输条件。解还必须属于谱隙中向相应无穷远衰减的分支。
具体而言，存在与 $k$ 有关的
$c_{\mathrm{dec},\pm},\alpha_\pm>0$，使对 $j\geq0$ 有
\begin{align}
  \|u_-\|_{H^1(\mathcal{C}_--jLe_x)}
  &\leq c_{\mathrm{dec},-}\ee^{-\alpha_-j}
    \|\varphi\|_{H^{1/2}_\beta(\Gamma_-)},
  &
  \|u_+\|_{H^1(\mathcal{C}_++jLe_x)}
  &\leq c_{\mathrm{dec},+}\ee^{-\alpha_+j}
    \|\varphi\|_{H^{1/2}_\beta(\Gamma_+)}.
  \label{eq:half-guide-decay}
\end{align}
这个条件比含糊的辐射条件更强，并与导模在 $B$ 中的有限能量以及所选投影谱隙一致。

谱隙条件本身并不保证问题~\eqref{eq:half-guide-dirichlet-problem} 唯一可解。记
$\Sigma_{D,\pm}(\beta)\subset G_\beta$ 为在 $\Gamma_\pm$ 施加齐次 Dirichlet 条件后，
相应半条带 $B_\pm$ 上仍存在非零衰减解的例外谱。Fliss 的连续模型在本质谱外除一个可数例外集外给出
半波导问题的适定性\cite[定理~4.1]{Fliss2013}；在当前分片常数模型中，这一例外集必须针对
包含传输条件的半波导算子定义，不能从连续折射率模型直接照搬。下文只考虑
\begin{align}
  k^2\in G_\beta\setminus
  \bigl(\Sigma_{D,-}(\beta)\cup\Sigma_{D,+}(\beta)\bigr),
  \label{eq:dtn-admissible-parameters}
\end{align}
并要求问题~\eqref{eq:half-guide-dirichlet-problem} 对每个 $\varphi$ 唯一可解。

在参数条件~\eqref{eq:dtn-admissible-parameters} 下，定义连续 DtN 算子
\begin{align}
  \Lambda_\pm(k):H^{1/2}_\beta(\Gamma_\pm)
  &\longrightarrow H^{-1/2}_\beta(\Gamma_\pm),
  &
  \Lambda_\pm(k)\varphi
  &:=\left.\partial_{\nu_\pm}u_\pm(\varphi;k)
    \right|_{\Gamma_\pm}.
  \label{eq:continuous-dtn}
\end{align}
其中法向导数按弱共法向迹理解。由于 $\nu_\pm$ 从中心指向 $B_\pm$，半条带自身外法向下的
导数等于 $-\Lambda_\pm(k)\varphi$。在 Dirichlet 例外参数处，单值算子
\eqref{eq:continuous-dtn} 可能失效；更稳健的连续对象是所有衰减半波导解的 Dirichlet--Neumann
Cauchy 数据所组成的线性关系。只有该关系到 Dirichlet 分量的投影可逆时，它才是
$\Lambda_\pm(k)$ 的图。本节始终处在式~\eqref{eq:dtn-admissible-parameters} 的单值情形。

利用这两个算子，可以把无界问题限制到 $\mathcal{C}_0$。中心场
$u_0\in H^1_\beta(\mathcal{C}_0)\setminus\{0\}$ 应满足
\begin{equation}
  \left\{
  \begin{aligned}
    -\Delta u_0&=k^2\rho u_0
    &&\text{在 }\mathcal{C}_0\setminus\Upsilon\text{ 的各均匀材料子区域内},\\
    u_0^+&=u_0^- ,
    &\partial_\nu u_0^+&=\partial_\nu u_0^-
    &&\text{在 }\Upsilon\cap\mathcal{C}_0,\\
    u_0|_{y=d/2}
    &=\ee^{\ii\beta d}u_0|_{y=-d/2},
    &\partial_yu_0|_{y=d/2}
    &=\ee^{\ii\beta d}\partial_yu_0|_{y=-d/2},\\
    \partial_{\nu_\pm}u_0
    &=\Lambda_\pm(k)\gamma_\pm u_0
    &&\text{在 }\Gamma_\pm.
  \end{aligned}
  \right.
  \label{eq:center-dtn-problem}
\end{equation}
最后一行是精确透明边界条件。对任意 $v\in H^1_\beta(\mathcal{C}_0)$ 分部积分时，顶、底
边界项因准周期条件相消，而 $\Gamma_\pm$ 的中心外法向恰为 $\nu_\pm$。代入
式~\eqref{eq:center-dtn-problem} 的最后一行，得到弱形式
\begin{align}
  \int_{\mathcal{C}_0}\nabla u_0\mathbin{\cdot}\nabla\overline v
  -k^2\int_{\mathcal{C}_0}\rho u_0\overline v
  -\sum_{\varsigma\in\{-,+\}}
   \left\langle
   \Lambda_\varsigma(k)\gamma_\varsigma u_0,
   \gamma_\varsigma v
   \right\rangle_{\Gamma_\varsigma}
  &=0.
  \label{eq:center-dtn-weak-form}
\end{align}
两个 DtN 项的负号都来自中心区域的分部积分与
式~\eqref{eq:artificial-interfaces} 的法向约定。由此定义连续非线性算子函数
\begin{align}
  F_\beta(k):H^1_\beta(\mathcal{C}_0)
  &\longrightarrow H^1_\beta(\mathcal{C}_0)^*,
  \notag\\
  \langle F_\beta(k)u,v\rangle
  &:=\int_{\mathcal{C}_0}\nabla u\mathbin{\cdot}\nabla\overline v
  -k^2\int_{\mathcal{C}_0}\rho u\overline v
  -\sum_{\varsigma\in\{-,+\}}
   \left\langle
   \Lambda_\varsigma(k)\gamma_\varsigma u,
   \gamma_\varsigma v
   \right\rangle_{\Gamma_\varsigma}.
  \label{eq:continuous-center-operator-function}
\end{align}
它关于 $k$ 非线性，因为 $\Lambda_\pm(k)$ 本身依赖于 $k$。

\begin{proposition}[无界导模问题与中心 DtN 问题的等价性]
  \label{prop:continuous-dtn-equivalence}
  固定满足假设~\ref{ass:projected-gap} 的 $\beta$，并令 $k>0$ 满足
  式~\eqref{eq:dtn-admissible-parameters}，且两个半波导 Dirichlet 问题均唯一可解。
  则 $k^2$ 是原无界 $\beta$-表述~\eqref{eq:gap-restricted-beta-formulation} 的导模特征值，
  当且仅当
  \begin{align}
    \ker F_\beta(k)&\ne\{0\}.
    \label{eq:dtn-kernel-criterion}
  \end{align}
  此外，全局导模到 $\mathcal{C}_0$ 的限制给出两个解空间之间的线性双射。
\end{proposition}

\begin{proof}
  设 $u$ 是原无界问题的非零导模。它在 $B_\pm$ 上的限制满足
  问题~\eqref{eq:half-guide-dirichlet-problem}，其 Dirichlet 数据为
  $\varphi_\pm=\gamma_\pm(u|_{\mathcal{C}_0})$。由半波导问题的唯一性，这两个限制正是定义
  $\Lambda_\pm(k)$ 的衰减解，故沿同一法向 $\nu_\pm$ 的法向导数与中心限制相等，并满足
  式~\eqref{eq:center-dtn-problem} 的透明边界条件。分部积分于是给出
  $F_\beta(k)(u|_{\mathcal{C}_0})=0$。若中心限制为零，则其界面迹为零，半波导唯一性又迫使
  两侧限制均为零，这与 $u\ne0$ 矛盾；故 $\ker F_\beta(k)\ne\{0\}$。

  反之，取 $u_0\in\ker F_\beta(k)\setminus\{0\}$。以
  $\varphi_\pm=\gamma_\pm u_0$ 分别求解两个唯一的衰减半波导 Dirichlet 问题，并得到
  $u_-$、$u_+$。弱式~\eqref{eq:continuous-center-operator-function} 先由内部试验函数给出中心区域内的
  分片 Helmholtz 方程和传输条件，再由界面迹的任意性给出
  $\partial_{\nu_\pm}u_0=\Lambda_\pm(k)\gamma_\pm u_0$。因此三部分在
  $\Gamma_\pm$ 上具有相同的 Dirichlet 迹和按同一法向计算的法向导数。将
  $u_-$、$u_0$、$u_+$ 依次拼接，便得到 $B$ 上满足全部内部传输条件和 $y$ 方向
  $\beta$-准周期条件的弱解。式~\eqref{eq:half-guide-decay} 保证拼接后的场属于
  $H^1_\beta(B)$，所以它是原问题的导模；其中心部分非零，故全局场非零。两种构造互为逆，
  从而还得到所述线性双射。
\end{proof}

命题~\ref{prop:continuous-dtn-equivalence} 是针对当前分片常数折射率与传输条件写出的连续
版本；其限制与延拓结构对应于 Fliss 的有界域 DtN 等价定理\cite[定理~4.5]{Fliss2013}。
下一节将构造连续 DtN 算子的数值近似，并据此形成离散非线性特征值问题。


\section{导模与特征值的数值计算}
\label{sec:numerics}

第~\ref{sec:dtn}~节中的 $\Lambda_\pm(k)$ 由半无限偏微分方程定义，是连续迹空间之间的
精确算子。本节依次离散周期胞元、界面迹和中心胞元，所得 Cauchy 数据空间、离散 DtN
以及最终有限矩阵都只是这一连续约化的数值实现，不能反过来定义 $\Lambda_\pm(k)$ 或
$F_\beta(k)$。以下仍固定 $\beta$，并只考虑
$k\in(\sqrt{g_-},\sqrt{g_+})$；材料系数继续采用分片常数模型，边界积分方程在每条
材料界面上同时施加场和法向导数的传输条件。

首先截断人工界面上的 Fourier 展开。令
\begin{align}
  \beta_m&=\beta+\frac{2\pi m}{d},
  &\psi_m(y)&=d^{-1/2}\ee^{\ii\beta_m y},
  &m&=-M,\ldots,M,
  &K&=2M+1.
  \label{eq:port-fourier-basis}
\end{align}
这些基函数自动满足 $y$ 方向的 $\beta$-准周期条件。固定截断阶数 $M$，并将截断系数列
$c$ 按零延拓理解，即 $|m|>M$ 时 $c_m=0$。
Dirichlet 与 Neumann 系数分别用 $H^{1/2}$ 和 $H^{-1/2}$ 的自然权重衡量：
\begin{align}
  \lVert c\rVert_{s,\beta}
  &:=\left\lVert
    \sum_{m\in\mathbb Z}c_m\psi_m
  \right\rVert_{H^s_\beta(\Gamma_\pm)}
  :=\left(
    \sum_{m\in\mathbb Z}
    (1+|\beta_m|^2)^s|c_m|^2
  \right)^{1/2},
  \qquad s\in\left\{\frac12,-\frac12\right\}.
  \label{eq:port-fourier-norms}
\end{align}
因此 Dirichlet 系数 $d$ 使用 $\lVert d\rVert_{1/2,\beta}$，Neumann 系数 $n$
使用 $\lVert n\rVert_{-1/2,\beta}$；$\beta=0$ 时可简写为 $\lVert\cdot\rVert_s$。两类系数的底层
空间均为 $\CC^K$，但分别带有上述 $s=1/2$ 与 $s=-1/2$ 范数，对偶配对取为
$d^*n$。改变 $M$ 时，Dirichlet 系数的零填充与 Neumann 系数的对偶限制必须配套。

对左右基本周期胞元 $\mathcal C_-$ 和 $\mathcal C_+$，在
$\partial\Omega_-$、$\partial\Omega_+$ 上以层势表示各均匀材料子区域内的场，并用边界
积分方程施加 transmission conditions；胞元顶、底边界保持 $\beta$-准周期匹配，两个竖直端口
则以式~\eqref{eq:port-fourier-basis} 的入射和出射振幅作为数据。采用步长统记为 $h$ 的
高阶 Nystr\"om 离散后，可逐列求解得到一个周期胞元的反射与透射块
\cite{LuLu2014}：
\begin{align}
  \begin{bmatrix}b_{\mathrm L}\\ a_{\mathrm R}\end{bmatrix}
  &=
  \begin{bmatrix}
    R_{\mathrm L}&T_{\mathrm{RL}}\\
    T_{\mathrm{LR}}&R_{\mathrm R}
  \end{bmatrix}
  \begin{bmatrix}a_{\mathrm L}\\ b_{\mathrm R}\end{bmatrix},
  \qquad
  R_{\mathrm L},T_{\mathrm{RL}},T_{\mathrm{LR}},R_{\mathrm R}
  \in\CC^{K\times K}.
  \label{eq:one-cell-scattering}
\end{align}
这里 $\mathrm L$、$\mathrm R$ 标记同一周期胞元的左、右端口；它们不表示半无限条带
$B_-$、$B_+$。
式~\eqref{eq:one-cell-scattering} 是胞元边界积分离散与端口截断共同产生的有限维传播关系。
从此处起固定离散层 $(h,M)$。除了由第~\ref{sec:dtn}~节继承的连续对象
$\Lambda_\pm(k)$ 和 $F_\beta(k)$ 外，后文的 $J(k)$、$\mathsf F_\pm(k)$、
$\Pi_{\mathrm D,\pm}(k)$、$\Pi_{\mathrm N,\pm}(k)$、$\mathsf V_\pm(k)$、$A(k)$
与 $A_{\mathrm{aug}}(k)$ 都表示固定 $(h,M)$ 后的有限矩阵或有限维对象，不再逐项重复
标记 $(h,M)$。只有确实比较不同离散层时才使用完整上标 ${}^{(h,M)}$。

记 $z_{\mathrm L}=(a_{\mathrm L},b_{\mathrm L})^{\mathsf T}$、
$z_{\mathrm R}=(a_{\mathrm R},b_{\mathrm R})^{\mathsf T}$。将
式~\eqref{eq:one-cell-scattering} 改写为
\begin{align}
  A_{\mathrm{sc}}z_{\mathrm L}&=B_{\mathrm{sc}}z_{\mathrm R},
  &
  A_{\mathrm{sc}}&=
  \begin{bmatrix}
    -R_{\mathrm L}&I\\ T_{\mathrm{LR}}&0
  \end{bmatrix},
  &
  B_{\mathrm{sc}}&=
  \begin{bmatrix}
    0&T_{\mathrm{RL}}\\ I&-R_{\mathrm R}
  \end{bmatrix}.
  \label{eq:one-cell-matrix-pair}
\end{align}
相邻胞元状态满足 $z_{\mathrm R}=\lambda z_{\mathrm L}$，故胞元传播由广义特征值问题
\begin{align}
  A_{\mathrm{sc}}z&=\lambda B_{\mathrm{sc}}z
  \label{eq:one-cell-generalized-eigenproblem}
\end{align}
决定。$\lambda$ 是沿 $x$ 方向平移一个周期 $L$ 的 Floquet 乘子，不是本文要求的特征值
$k^2$。数值上对矩阵对 $(A_{\mathrm{sc}},B_{\mathrm{sc}})$ 作有序 QZ 分解：原问题选取
朝 $x\to+\infty$ 衰减的挠曲子空间，反向关系
$B_{\mathrm{sc}}z_{\mathrm R}=\mu A_{\mathrm{sc}}z_{\mathrm R}$ 则独立选取朝
$x\to-\infty$ 衰减的子空间。
在实谱隙中的参考点，所选两簇各应有 $K$ 维，并分别由 $|\lambda|<1$ 和
$|\mu|<1$ 识别。若矩阵对不正则、出现无法分离的单位圆谱簇，或衰减子空间的维数与参考点
不一致，则该参数处的半波导离散不可用，不应通过逐点改选特征向量掩盖这一失效。

令 $Z_+(k)\in\CC^{2K\times K}$ 为正向问题
$A_{\mathrm{sc}}z=\lambda B_{\mathrm{sc}}z$ 中向 $x\to+\infty$ 衰减的广义 Schur 子空间的一组基，
$Z_-(k)\in\CC^{2K\times K}$ 为反向问题
$B_{\mathrm{sc}}z=\mu A_{\mathrm{sc}}z$ 中向 $x\to-\infty$ 衰减的广义 Schur 子空间的一组基。
记材料界面外与端口所在的均匀背景系数为 $\rho_{\mathrm e}>0$。沿用 Rayleigh 通道的分支，令
\begin{align}
  \gamma_m(k)^2&=k^2\rho_{\mathrm e}-\beta_m^2,
  &\Im\gamma_m(k)&\geq0,
  &\gamma_m(k)&\geq0\quad\text{当 $\gamma_m(k)$ 为实数时},
  \notag\\
  \boldsymbol{\gamma}(k)
  &:=\operatorname{diag}\bigl(
    \gamma_{-M}(k),\ldots,\gamma_M(k)
  \bigr).
  \label{eq:rayleigh-gamma-matrix}
\end{align}
按照 $\nu_-=-e_x$、$\nu_+=e_x$ 将这两组振幅基转换为人工边界上的离散 Cauchy 数据基：
\begin{align}
  \mathsf V_-(k)
  &:={}
  \begin{bmatrix}
    I&I\\
    -\ii\boldsymbol{\gamma}(k)&\ii\boldsymbol{\gamma}(k)
  \end{bmatrix}Z_-(k),
  &
  \mathsf V_+(k)
  &:={}
  \begin{bmatrix}
    I&I\\
    \ii\boldsymbol{\gamma}(k)&-\ii\boldsymbol{\gamma}(k)
  \end{bmatrix}Z_+(k).
  \label{eq:qz-to-cauchy-basis}
\end{align}
再按行分块
\begin{align}
  \mathsf V_\pm(k)
  &:={}
  \begin{bmatrix}
    \mathsf V_{\mathrm D,\pm}(k)\\
    \mathsf V_{\mathrm N,\pm}(k)
  \end{bmatrix},
  &
  \mathsf V_{\mathrm D,\pm}(k),
  \mathsf V_{\mathrm N,\pm}(k)
  &\in\CC^{K\times K},
  \notag\\
  \mathcal G_\pm(k)
  &:=\operatorname{range}\mathsf V_\pm(k)
  \subset\CC^K\times\CC^K.
  \label{eq:discrete-cauchy-space}
\end{align}
$\mathsf V_\pm(k)$ 的每一列是一个衰减半波导解在 $\Gamma_\pm$ 上的离散 Cauchy
数据对，其列构成 $\mathcal G_\pm(k)$ 的一组基。上、下两个 $\CC^K$ 分别带有
$\lVert\cdot\rVert_{1/2,\beta}$ 和 $\lVert\cdot\rVert_{-1/2,\beta}$ 范数。任意允许的 Cauchy 数据可写为
\begin{align}
  d_\pm&=\mathsf V_{\mathrm D,\pm}(k)c_\pm,
  &n_\pm&=\mathsf V_{\mathrm N,\pm}(k)c_\pm,
  &c_\pm&\in\CC^K.
  \label{eq:discrete-cauchy-relation}
\end{align}
更换有序 QZ 基只会使 $\mathsf V_\pm(k)$ 右乘一个可逆矩阵，因而不改变其列空间，也不改变
下述坐标安全时形成的离散 DtN。

只有当 $\mathsf V_{\mathrm D,\pm}(k)$ 在式~\eqref{eq:port-fourier-norms} 的加权坐标中可安全求解，
且 Dirichlet 投影与 $\mathcal G_\pm(k)$ 安全横截时，才定义离散 DtN 矩阵
\begin{align}
  \widetilde{\Lambda}_\pm(k)
  &:=\mathsf V_{\mathrm N,\pm}(k)
    \mathsf V_{\mathrm D,\pm}(k)^{-1}:
  \CC^K\longrightarrow\CC^K.
  \label{eq:discrete-dtn-matrix}
\end{align}
定义域与值域上的 $\CC^K$ 分别带有 $s=1/2$ 与 $s=-1/2$ 范数。实现中用右线性
求解计算式~\eqref{eq:discrete-dtn-matrix}，不显式求逆。除矩阵条件数外，还需检查 Dirichlet 投影
相对于整个 Cauchy 数据空间的绝对横截余量；一个条件数不大的 $\mathsf V_{\mathrm D,\pm}(k)$
仍可能对应近乎竖直的 Cauchy 数据空间。若任一 Dirichlet 坐标不安全，则保留
式~\eqref{eq:discrete-cauchy-relation} 中的完整 Cauchy 关系，而不强行形成 $\widetilde{\Lambda}_\pm(k)$。
这一区分也覆盖连续 Dirichlet 例外谱附近的数值情形；波浪号在此仅表示它是连续
$\Lambda_\pm(k)$ 的数值近似，不是所有离散矩阵的通用标记。

中心胞元 $\mathcal C_0$ 使用相同的分片 Helmholtz 层势和 transmission conditions 离散。
为说明离散未知量的来源，先给出相应的连续表示。若
$\Omega_0=\bigcup_{\ell=1}^{N_0}\Omega_{0,\ell}$，在每个 $\partial\Omega_{0,\ell}$ 上以 $\nu$ 表示由介质柱
内部指向外部的单位法向。对
$\mathbf z\notin\partial\Omega_{0,\ell}$ 定义
\begin{align*}
  \mathcal S_{\mathrm{QP},\ell}^{(\kappa)}[\sigma](\mathbf z)
  &:=\int_{\partial\Omega_{0,\ell}}
    G_{\mathrm{QP}}^{(\kappa)}(\mathbf z,\mathbf z')
    \sigma(\mathbf z')\,\dd s(\mathbf z'),\\
  \mathcal D_{\mathrm{QP},\ell}^{(\kappa)}[\tau](\mathbf z)
  &:=\int_{\partial\Omega_{0,\ell}}
    \partial_{\nu(\mathbf z')}G_{\mathrm{QP}}^{(\kappa)}(\mathbf z,\mathbf z')
    \tau(\mathbf z')\,\dd s(\mathbf z').
\end{align*}
这里 $G_{\mathrm{QP}}^{(\kappa)}$ 是沿 $y$ 方向具有相位
$\ee^{\ii\beta d}$ 的准周期 Green 函数；以自由空间 Green 函数
$G^{(\kappa)}$ 替换 $G_{\mathrm{QP}}^{(\kappa)}$，得到
$\mathcal S_{\ell}^{(\kappa)}$ 和 $\mathcal D_{\ell}^{(\kappa)}$。

\begin{theorem}[中心胞元的层势与均匀背景场表示]
  \label{thm:center-cell-representation}
  假设各 $\partial\Omega_{0,\ell}$ 两两分离、具有 $C^2$ 正则性并严格包含在
  $\mathcal C_0$ 内，$\mathcal C_0\setminus\overline{\Omega_0}$ 中的系数为
  $\rho_{\mathrm e}$，而 $\Omega_{0,\ell}$ 中的系数为正的常数
  $\rho_{\mathrm i,\ell}$。令
  $\kappa_{\mathrm e}=k\sqrt{\rho_{\mathrm e}}$、
  $\kappa_{\mathrm i,\ell}=k\sqrt{\rho_{\mathrm i,\ell}}$，并假设
  式~\eqref{eq:rayleigh-gamma-matrix} 的既定分支满足
  $\gamma_m(k)\neq0$，$m\in\mathbb Z$。再假设对每个材料界面，交换内、外波数后所得的
  全平面 outgoing transmission problem 对任意界面跳跃数据唯一可解。

  若 $u\in H^1_\beta(\mathcal C_0)$ 在各均匀材料子区域内属于 $H^2$，逐片满足
  Helmholtz 方程并在 $\partial\Omega_0$ 上满足 transmission conditions，而在
  $\Gamma_-$、$\Gamma_+$ 上不预先施加边界条件，则存在密度
  $\tau_\ell\in H^{1/2}(\partial\Omega_{0,\ell})$、
  $\sigma_\ell\in H^{-1/2}(\partial\Omega_{0,\ell})$ 以及唯一的 Rayleigh 系数列
  $(a_{0,m}^-)_{m\in\mathbb Z}$、$(b_{0,m}^+)_{m\in\mathbb Z}$，使得
  \begin{equation*}
    u(\mathbf z)=
    \begin{cases}
      \displaystyle
      \sum_{\ell=1}^{N_0}\left(
        \mathcal S_{\mathrm{QP},\ell}^{(\kappa_{\mathrm e})}[\sigma_\ell]
        +\mathcal D_{\mathrm{QP},\ell}^{(\kappa_{\mathrm e})}[\tau_\ell]
      \right)(\mathbf z)
      +u_{\mathrm{hom}}(\mathbf z),
      &\mathbf z\in\mathcal C_0\setminus\overline{\Omega_0},\\[2mm]
      \displaystyle
      \left(
        \mathcal S_{\ell}^{(\kappa_{\mathrm i,\ell})}[\sigma_\ell]
        +\mathcal D_{\ell}^{(\kappa_{\mathrm i,\ell})}[\tau_\ell]
      \right)(\mathbf z),
      &\mathbf z\in\Omega_{0,\ell},
      \quad 1\leq\ell\leq N_0,
    \end{cases}
  \end{equation*}
  其中均匀背景场为
  \begin{equation*}
    u_{\mathrm{hom}}(x,y)
    =\sum_{m\in\mathbb Z}a_{0,m}^-
      \ee^{\ii\gamma_m(k)(x+a)}\psi_m(y)
     +\sum_{m\in\mathbb Z}b_{0,m}^+
      \ee^{-\ii\gamma_m(k)(x-a)}\psi_m(y),
  \end{equation*}
  两个级数按 $H^1(\mathcal C_0)$ 意义收敛。换言之，中心胞元外部场由准周期层势贡献与
  homogeneous background field 之和表示，内部场则由相应介质柱边界上的自由空间层势表示。
  若将外部准周期层势与内部自由空间层势交换适用区域后形成的齐次互补层势系统只有零解，
  则上述层势密度也唯一。
\end{theorem}

定理~\ref{thm:center-cell-representation} 的证明见
附录~\ref{app:center-cell-representation}。它也说明准周期层势本身不能表示全部中心胞元场：
即使 $\Omega_0=\varnothing$，层势贡献消失，$u_{\mathrm{hom}}$ 仍可能非零。
对定理中的密度作 Nystr\"om 离散并将两组 Rayleigh 级数截断到 $-M\leq m\leq M$ 后，
仍以 $\tau$、$\sigma$ 表示密度系数，并令 $n$ 为材料界面
密度未知量的总数；再令 $a_0^-,b_0^+\in\CC^K$ 分别表示锚定在
$\Gamma_-$、$\Gamma_+$ 的 homogeneous Rayleigh 振幅，并置
\begin{align}
  \eta&=
  \begin{bmatrix}\tau\\-\sigma\end{bmatrix}\in\CC^n,
  &\xi&=
  \begin{bmatrix}\eta\\a_0^-\\b_0^+\end{bmatrix}
  \in\CC^{n+2K}.
  \label{eq:center-discrete-unknown}
\end{align}
这里 $\eta$ 只收集 $\partial\Omega_0$ 上的层势密度系数，而 $\xi$ 是中心胞元离散系统的完整
系数向量。因而 $a_0^-$ 和 $b_0^+$ 必须保留；当中心胞元没有材料界面时 $n=0$，
只有密度及相应的材料界面方程消失。固定 $h$ 和 $M$ 后始终使用 $\xi$；只有比较不同离散层时
才写作 $\xi^{(h,M)}$。

以一个背景材料与介质柱的界面为例，记界面外、内两侧的分片常数系数分别为
$\rho_{\mathrm e}$、$\rho_{\mathrm i}$，并令
\begin{align}
  \kappa_{\mathrm e}&=k\sqrt{\rho_{\mathrm e}},
  &\kappa_{\mathrm i}&=k\sqrt{\rho_{\mathrm i}}.
  \label{eq:center-material-wavenumbers}
\end{align}
以下 $S,D,D^*,T$ 均表示当前 $h$ 网格上的 Nystr\"om 矩阵；带 $\mathrm{QP}$ 下标的版本
使用 $y$ 方向准周期 Green 函数，不带该下标的版本使用自由空间 Green 函数，$T$ 是
hypersingular operator。令 $Q_{\mathrm D,-},Q_{\mathrm N,-},Q_{\mathrm D,+},Q_{\mathrm N,+}$
表示两组 homogeneous Rayleigh 场在材料界面上的 Dirichlet 和 Neumann 迹矩阵，其条目将在
式~\eqref{eq:center-homogeneous-rayleigh-traces} 中给出。于是作用于完整向量 $\xi$ 的材料界面
jump/mismatch 矩阵为
\begin{align}
  J(k)
  &=
  \begin{bmatrix}
    I+D_{\mathrm{QP}}^{(\kappa_{\mathrm e})}
      -D^{(\kappa_{\mathrm i})}
    &-S_{\mathrm{QP}}^{(\kappa_{\mathrm e})}
      +S^{(\kappa_{\mathrm i})}
    &Q_{\mathrm D,-}&Q_{\mathrm D,+}\\
    T_{\mathrm{QP}}^{(\kappa_{\mathrm e})}
      -T^{(\kappa_{\mathrm i})}
    &I-D_{\mathrm{QP}}^{(\kappa_{\mathrm e})*}
      +D^{(\kappa_{\mathrm i})*}
    &Q_{\mathrm N,-}&Q_{\mathrm N,+}
  \end{bmatrix}
  \in\CC^{n\times(n+2K)}.
  \label{eq:center-interface-jump-matrix}
\end{align}
若 $\Omega_0$ 含有多根介质柱，则这些块在全部材料界面上整体组装：对角块包含 jump terms，
非对角块表示不同界面之间的层势相互作用。因此 $J(k)\xi=0$ 同时施加函数值 jump 和
法向导数 jump 为零。

为明确 homogeneous Rayleigh 场在材料界面上的两组列，取 Nystr\"om 节点与介质柱外法向
\begin{align*}
  z_j&=(x_j,y_j)\in\partial\Omega_0,
  &\nu_j&=(\nu_{x,j},\nu_{y,j}),
\end{align*}
并沿用式~\eqref{eq:rayleigh-gamma-matrix} 的 $\gamma_m(k)$ 分支。于是对 $-M\leq m\leq M$，
\begin{align}
  (Q_{\mathrm D,-})_{jm}
  &=\ee^{\ii\gamma_m(k)(x_j+a)}\psi_m(y_j),
  &
  (Q_{\mathrm N,-})_{jm}
  &=\ii(\gamma_m(k)\nu_{x,j}+\beta_m\nu_{y,j})
    \ee^{\ii\gamma_m(k)(x_j+a)}\psi_m(y_j),
  \notag\\
  (Q_{\mathrm D,+})_{jm}
  &=\ee^{-\ii\gamma_m(k)(x_j-a)}\psi_m(y_j),
  &
  (Q_{\mathrm N,+})_{jm}
  &=\ii(-\gamma_m(k)\nu_{x,j}+\beta_m\nu_{y,j})
    \ee^{-\ii\gamma_m(k)(x_j-a)}\psi_m(y_j).
  \label{eq:center-homogeneous-rayleigh-traces}
\end{align}

下面直接由准周期 Green 函数的 Rayleigh 谱展开提取中心层势在人工边界上的
outgoing Dirichlet 系数。以上述 Nystr\"om 节点 $z_j=(x_j,y_j)$、源法向
$\nu_j=(\nu_{x,j},\nu_{y,j})$ 和求积权重 $w_j$ 为例，对 $-M\leq m\leq M$ 令
\begin{align}
  g_-(m,j)
  &:=\frac{\ii}{2\sqrt{d}\,\gamma_m(k)}
  \exp\!\left(
    -\ii\beta_m y_j+\ii\gamma_m(k)(x_j+a)
  \right),
  \notag\\
  g_+(m,j)
  &:=\frac{\ii}{2\sqrt{d}\,\gamma_m(k)}
  \exp\!\left(
    -\ii\beta_m y_j-\ii\gamma_m(k)(x_j-a)
  \right),
  \label{eq:center-farfield-kernels}\\
  \partial_{\nu_j}g_-(m,j)
  &=\ii\bigl(\gamma_m(k)\nu_{x,j}-\beta_m\nu_{y,j}\bigr)g_-(m,j),
  \notag\\
  \partial_{\nu_j}g_+(m,j)
  &=-\ii\bigl(\gamma_m(k)\nu_{x,j}+\beta_m\nu_{y,j}\bigr)g_+(m,j).
  \label{eq:center-farfield-normal-factors}
\end{align}
按 $\eta=[\tau;-\sigma]^{\mathsf T}$ 的列顺序，以模态 $m$ 为行、边界节点 $j$ 为列组装
\begin{align}
  \mathsf F_-(k)
  &:={}
  \left[
    \bigl(w_j\partial_{\nu_j}g_-(m,j)\bigr)_{m,j}
    \quad
    \bigl(-w_jg_-(m,j)\bigr)_{m,j}
  \right],
  \notag\\
  \mathsf F_+(k)
  &:={}
  \left[
    \bigl(w_j\partial_{\nu_j}g_+(m,j)\bigr)_{m,j}
    \quad
    \bigl(-w_jg_+(m,j)\bigr)_{m,j}
  \right],
  \notag\\
  \mathsf F_-(k),\mathsf F_+(k)
  &\in\CC^{K\times n}.
  \label{eq:center-farfield-extractors}
\end{align}
因而 $\mathsf F_\pm(k)\eta$ 是中心材料界面层势在 $\Gamma_\pm$ 上的 outgoing Rayleigh
Dirichlet 系数，其 Neumann 系数由逐模法向导数乘子得到。$\mathsf F_\pm(k)$ 是根据谱展开与
Nystr\"om 求积直接组装的，不是先在人工边界上采样再作 FFT 或 DFT。若有多根介质柱，
则上述节点、权重与密度块在全部材料界面上整体组装。这一有限矩阵与连续非线性算子函数
$F_\beta(k)$ 的字体和定义均不同。

两组 homogeneous Rayleigh 振幅分别锚定在
$\Gamma_- = \{-a\}\times I_d$ 和 $\Gamma_+ = \{a\}\times I_d$。定义中心胞元传播矩阵
\begin{align}
  E_0(k)
  &:=\operatorname{diag}\left(
    \ee^{2a\ii\gamma_{-M}(k)},\ldots,
    \ee^{2a\ii\gamma_M(k)}
  \right).
  \label{eq:center-cell-propagation-matrix}
\end{align}
这里的传播长度是中心胞元宽度 $2a$，不是左右周期胞元的周期 $L$。于是中心场在
$\Gamma_\pm$ 上的 Dirichlet 与从 $\mathcal C_0$ 指向 $B_\pm$ 的 Neumann 系数由
\begin{align}
  \Pi_{\mathrm D,-}(k)
  &:={}
  \begin{bmatrix}
    \mathsf F_-(k)&I&E_0(k)
  \end{bmatrix},
  &
  \Pi_{\mathrm N,-}(k)
  &:={}
  \begin{bmatrix}
    \ii\boldsymbol{\gamma}(k)\mathsf F_-(k)
    &-\ii\boldsymbol{\gamma}(k)
    &\ii\boldsymbol{\gamma}(k)E_0(k)
  \end{bmatrix},
  \notag\\
  \Pi_{\mathrm D,+}(k)
  &:={}
  \begin{bmatrix}
    \mathsf F_+(k)&E_0(k)&I
  \end{bmatrix},
  &
  \Pi_{\mathrm N,+}(k)
  &:={}
  \begin{bmatrix}
    \ii\boldsymbol{\gamma}(k)\mathsf F_+(k)
    &\ii\boldsymbol{\gamma}(k)E_0(k)
    &-\ii\boldsymbol{\gamma}(k)
  \end{bmatrix},
  \label{eq:center-trace-extraction-matrices}\\
  \Pi_{\mathrm D,\pm}(k),
  \Pi_{\mathrm N,\pm}(k)
  &\in\CC^{K\times(n+2K)}.
  \notag
\end{align}
每个矩阵的第一列块来自中心材料界面的层势密度，第二、三列块分别来自锚定在
$\Gamma_-$、$\Gamma_+$ 的 $a_0^-$、$b_0^+$。Neumann 行块的符号由
$\nu_-=-e_x$、$\nu_+=e_x$ 直接推出。当中心胞元为空、$n=0$ 时，第一列块消失，上式退化为
$I$、$E_0(k)$ 与 $\boldsymbol{\gamma}(k)$ 的组装。

在两个 Dirichlet 坐标都安全时，中心离散非线性特征值问题写成矩阵函数方程
\begin{align}
  A(k)\xi
  &:={}
  \begin{bmatrix}
    J(k)\xi\\
    \Pi_{\mathrm N,-}(k)\xi
      -\widetilde{\Lambda}_-(k)\Pi_{\mathrm D,-}(k)\xi\\
    \Pi_{\mathrm N,+}(k)\xi
      -\widetilde{\Lambda}_+(k)\Pi_{\mathrm D,+}(k)\xi
  \end{bmatrix}
  =0.
  \label{eq:discrete-center-dtn-matrix}
\end{align}
其中 $A(k)\in\CC^{(n+2K)\times(n+2K)}$。第一块施加中心材料界面的
transmission conditions，后两块分别是左右人工边界上的离散 Neumann--DtN mismatch。
式~\eqref{eq:discrete-center-dtn-matrix} 中的两个负号离散了
式~\eqref{eq:center-dtn-weak-form} 的透明边界项；法向仍从 $\mathcal C_0$ 指向 $B_\pm$。
$A(k)$ 是由中心胞元边界积分表示和离散透明边界条件直接组装的有限矩阵函数，并非连续弱
算子 $F_\beta(k)$ 的直接矩阵离散；借助连续中心 DtN 问题与无界问题的等价性，它的零空间
近似 $\ker F_\beta(k)$。

若 Dirichlet 坐标失效，则用更原始的 Cauchy 数据关系组装
\begin{align}
  A_{\mathrm{aug}}(k)
  \begin{bmatrix}\xi\\c_-\\c_+\end{bmatrix}
  &:={}
  \begin{bmatrix}
    J(k)\xi\\
    \Pi_{\mathrm D,-}(k)\xi
      -\mathsf V_{\mathrm D,-}(k)c_-\\
    \Pi_{\mathrm N,-}(k)\xi
      -\mathsf V_{\mathrm N,-}(k)c_-\\
    \Pi_{\mathrm D,+}(k)\xi
      -\mathsf V_{\mathrm D,+}(k)c_+\\
    \Pi_{\mathrm N,+}(k)\xi
      -\mathsf V_{\mathrm N,+}(k)c_+
  \end{bmatrix}
  =0,
  &
  A_{\mathrm{aug}}(k)
  &\in\CC^{(n+4K)\times(n+4K)}.
  \label{eq:discrete-center-graph-matrix}
\end{align}
这是保留完整 Cauchy 数据坐标的增广匹配矩阵。当
$\mathsf V_{\mathrm D,-}(k)$、$\mathsf V_{\mathrm D,+}(k)$ 均可安全求解时，从
式~\eqref{eq:discrete-center-graph-matrix} 消去 $c_-$、$c_+$，恰好得到
式~\eqref{eq:discrete-center-dtn-matrix}。因此两式描述同一个离散匹配条件；前者保留完整
Cauchy 数据关系，后者只是它在安全 Dirichlet 坐标下的降维形式。若无特别说明，
$A(k)$ 始终指式~\eqref{eq:discrete-center-dtn-matrix} 的默认离散 DtN 形式；只有
Dirichlet 坐标不安全时才改用 $A_{\mathrm{aug}}(k)$，不再强行形成离散 DtN。

候选定位只在所选实谱隙内进行。在默认离散 DtN 形式中，以固定的可逆行列权重
$W_r(k)$、$W_c(k)$ 平衡 Dirichlet、Neumann 与边界积分未知量的自然尺度，可采用
\begin{align}
  s(k)
  &:=\frac{\sigma_{\min}\!\left(
    W_r(k)A(k)W_c(k)\right)}
  {\sigma_{\max}\!\left(
    W_r(k)A(k)W_c(k)\right)}
  \label{eq:discrete-candidate-score}
\end{align}
作为无量纲定位量，先作实轴粗扫，再在孤立低谷附近使用固定规则逐层加密。若采用完整
Cauchy 数据关系，则在同一定义中把 $A(k)$ 换成 $A_{\mathrm{aug}}(k)$，并选取尺寸
相容的行列权重。权重不改变矩阵
的精确核，但会改变奇异值的尺度，因此候选必须同时通过原始矩阵残量、胞元边界积分残量、
QZ 子空间分离、Cauchy 数据匹配以及相邻离散层模式一致性检查。评分低谷只能产生离散特征值
候选；在尚未给出连续逼近与误差控制以前，不能把它直接称为原算子 $A(\beta)$ 的特征值。

在候选 $\widehat k$ 处，取相应矩阵函数的最小右奇异方向。其中心分量 $\xi$ 给出边界密度和
端口振幅，从而通过中心胞元层势重构 $u_{0,h,M}$。在 Cauchy 数据形式中，$c_-$、$c_+$
已经随核向量得到；在安全 DtN 形式中，则由
$\mathsf V_{\mathrm D,\pm}(k)c_\pm
=\Pi_{\mathrm D,\pm}(k)\xi$ 求得半波导坐标，并且相应的 Neumann 数据满足
$\mathsf V_{\mathrm N,\pm}(k)c_\pm
=\Pi_{\mathrm N,\pm}(k)\xi$。随后利用所选 QZ
挠曲子空间和单胞传播关系，在
$\mathcal C_--jLe_x$ 与 $\mathcal C_++jLe_x$，$j\geq0$，逐胞元重构向左右无穷远衰减的场。
三部分在 $\Gamma_\pm$ 上按 Dirichlet 和中心外法向 Neumann 数据拼接；各材料界面仍由原边界
积分表示满足 transmission conditions。最后按
\begin{align}
  w_{h,M}(x,y+jd)
  &=\ee^{\ii\beta jd}u_{h,M}(x,y),
  \qquad y\in I_d,\qquad j\in\mathbb Z,
  \label{eq:numerical-guided-mode-extension}
\end{align}
恢复全平面上的准周期导模，并以固定范数归一化。重构结果还应检查端口匹配残量、材料界面
传输残量和向两侧逐胞元衰减；这些检查用于判别数值伪影，不在本节给出未经计算的数值结论。

下一节将以这一离散候选和重构场为输入，讨论后验误差估计及其可靠性条件；本节不预先给出
误差上界。


\section{基于预解式的谱距离估计}
\label{sec:error}

第~\ref{sec:numerics}~节的有限矩阵函数只能定位候选并生成一个分片重构场。小矩阵残量
本身不是连续算子的残量，更不能直接证明物理谱点的存在。本节固定同一候选、同一分片场及其全部
材料和几何输入，直接评价该场在真实界面上的连续迹跳跃；随后用正移位预解式的全界面迹
估计，把这些跳跃转化为到 $\sigma(A(\beta))$ 的距离上界。这个论证不先把分片场提升到
全局 $H^1$ 空间，也不要求数值求解任何预解式方程。

\subsection{重构场与界面跳跃}

固定第~\ref{sec:numerics}~节得到的 $\widehat k>0$，并在本节始终使用
\begin{align}
  \mu&:=\widehat k^2.
  \label{eq:resolvent-candidate-mu}
\end{align}
在 $\widehat k$ 处固定第~\ref{sec:numerics}~节候选计算交付的中心系数 $\xi$、左右
Cauchy 坐标 $c_-$、$c_+$、单胞边界密度、端口截断、连续核评价规则和完整单胞传播关系，
以及由这些数据确定的分片连续核场 $u_{h,M}$。以下称 $u_{h,M}$ 为原始场。它不是重新
离散后的场，也不随求积精度的提高而改变。近核向量只确定到非零复数倍，
因此先固定第~\ref{sec:numerics}~节使用的非零归一化与相位约定；后面的 $U$、界面缺陷和
全部单侧界都必须由这个同一尺度的场计算。

这种固定数据的方式适用于一般中心缺陷，并不限定密度的具体存储方式。设候选计算在均匀参数网格上
交付偶数个密度样本 $\{\eta_r\}_{r=0}^{N_\eta-1}$，其中 $N_\eta$ 是密度样本数。
固定这些样本，并按其存储位型将其解释为精确输入，再由数学离散 Fourier 变换定义有限三角密度：
\begin{align}
  \eta_m
  &:=\frac1{N_\eta}\sum_{r=0}^{N_\eta-1}\eta_r
    \exp\!\left(-\frac{2\pi\ii mr}{N_\eta}\right),
  \qquad
  -\frac{N_\eta}{2}\leq m\leq\frac{N_\eta}{2}-1.
  \label{eq:exact-density-dft}
\end{align}
该有限密度作用于使用同一个精确 $\mu$ 的完整数学准周期 Green 核，并与固定的有限
Rayleigh 和共同定义 $u_{h,M}$。可靠上界计算中的核展开截断、局部级数截断和曲线求积
只用于评价这个既定场的实际连续迹跳跃；它们的截断误差与舍入误差进入相应球上界，而不
重新定义另一个原始场。数值 FFT 也只可作为式~\eqref{eq:exact-density-dft} 的经验证实现，
不能独立定义密度系数。

原始场在每个均匀材料子区域内由固定密度和指定的精确 Helmholtz 核定义，并满足
\begin{align}
  (-\Delta-\mu\rho)u_{h,M}&=0
  \qquad\text{逐片成立}.
  \label{eq:raw-piecewise-pde}
\end{align}
它也精确满足 $y$ 方向的 $\beta$-准周期条件。本文假定
$u_{h,M}\in\mathcal H$，且它在每个分片直至相应界面局部属于 $H^2$。这些条件使
逐片 Green 恒等式和普通法向迹有意义，但不假定
$u_{h,M}\in H^1_\beta(B)$：有限边界积分匹配和端口截断仍可能在界面上留下小的
值跳跃和法向导数跳跃。

为避免重复计数，定义互不相交的界面族
\begin{align}
  \mathscr E&=\mathscr E_{\mathrm{mat}}\mathbin{\dot\cup}
  \mathscr E_{\mathrm{art}}.
  \label{eq:certification-interface-family}
\end{align}
$\mathscr E_{\mathrm{mat}}$ 包含 $\partial\Omega_0$（若
$\Omega_0\ne\varnothing$）以及 $\partial\Omega_-$、$\partial\Omega_+$ 的全部半无限周期
平移；$\mathscr E_{\mathrm{art}}$ 包含中心与首个周期胞元之间的 $\Gamma_-$、$\Gamma_+$，
以及左右半波导中所有相邻胞元的共享竖直界面。每条物理界面只出现一次。由于原始场的
准周期值与 $\partial_y$ 迹精确匹配，条带上下边界的 Green 项相消，不另计入
$\mathscr E$。

在每条 $\Gamma\in\mathscr E$ 上固定从侧 $1$ 指向侧 $2$ 的法向 $\nu_\Gamma$，并定义
\begin{align}
  a_\Gamma
  &:=u_{h,M}^{(2)}|_\Gamma
    -u_{h,M}^{(1)}|_\Gamma,
  &j_\Gamma
  &:=\partial_{\nu_\Gamma}u_{h,M}^{(1)}|_\Gamma
    -\partial_{\nu_\Gamma}u_{h,M}^{(2)}|_\Gamma.
  \label{eq:raw-interface-jumps}
\end{align}
在材料界面上，侧 $1$ 是介质柱内部、侧 $2$ 是外部，$\nu_\Gamma$ 由内指向外；因此
$a_\Gamma=u^{\mathrm{out}}-u^{\mathrm{in}}$，而
$j_\Gamma=\partial_{\nu_\Gamma}u^{\mathrm{in}}
-\partial_{\nu_\Gamma}u^{\mathrm{out}}$。在人工界面上，侧 $1$ 是较靠近中心的胞元，侧
$2$ 是较远的胞元，法向朝相应无穷远；在 $\Gamma_\pm$ 上这正是
$\nu_-=-e_x$、$\nu_+=e_x$。若改用共享界面两侧各自的外法向，$j_\Gamma$ 等于两侧
外法向导数之和。这个定向使后面的 Green 恒等式中值跳跃项和法向跳跃项都带正号。

选定参考长度 $\ell_*>0$，定义实际界面缺陷
\begin{align}
  \varepsilon_{\partial,\ell_*}^2
  &:={}
  \sum_{\Gamma\in\mathscr E}
  \left(
    \lVert j_\Gamma\rVert_{L^2(\Gamma)}^2
    +\ell_*^{-2}\lVert a_\Gamma\rVert_{L^2(\Gamma)}^2
  \right),
  \label{eq:boundary-defect-norm}\\
  U&:=\lVert u_{h,M}\rVert_{\mathcal H}.
  \label{eq:raw-field-mass}
\end{align}
两者使用同一个 $\mu$、同一个原始场和同一套材料参数。式
~\eqref{eq:boundary-defect-norm} 评价的是连续迹跳跃，不是 $J(\widehat k)\xi$、
$A(\widehat k)\xi$ 或端口矩阵不匹配量的范数。

无限半波导的求和仍保留完整传播矩阵。令 $P_\pm$ 表示把对应衰减 Cauchy 坐标向无穷远
推进一个周期的矩阵，则第 $j$ 个周期胞元的状态是 $P_\pm^jc_\pm$。一个胞元中的值和
法向跳跃作用连同 $L^2$ 权重形成正半定矩阵
$\mathsf G_{\partial,\pm}$ 后，半波导部分具有
\begin{align}
  \sum_{\varsigma\in\{-,+\}}
  \sum_{j=0}^{\infty}
  (P_\varsigma^jc_\varsigma)^*
  \mathsf G_{\partial,\varsigma}
  (P_\varsigma^jc_\varsigma)
  \label{eq:full-propagation-boundary-sum}
\end{align}
的结构；中心与首胞之间不能由平移生成的两条界面另行计算。每条共享界面只能归属于这两类
中的一类。式~\eqref{eq:full-propagation-boundary-sum} 使用完整 $P_\pm$，不把传播方向假定
为正交标量模态，也不删除非正规或 Jordan 耦合。材料跳跃可以直接在真实曲线上先作有符号
组合再求 $L^2$ 范数；下文所谓“非 Fourier 材料评价”仅指这一做法。端口场、人工竖直
界面以及源密度的构造仍然使用第~\ref{sec:numerics}~节的 Rayleigh/Fourier 表示。

\subsection{正移位预解式的全界面迹估计}

下面给出将式~\eqref{eq:boundary-defect-norm} 与全谱距离连接起来的几何常数。假定
$\mathscr E$ 局部有限，每条材料界面是简单闭合 $C^2$ 曲线，人工界面是准周期圆柱上的
竖直闭曲线。存在共同的法向邻域半宽 $\delta_{\mathrm{col}}>0$、曲率上界
$\kappa_*$ 和有限覆盖重数 $N_{\mathrm{ov}}$，使对每条界面的弧长参数
$\boldsymbol r_\Gamma(s)$，映射
\begin{align}
  \boldsymbol\Phi_\Gamma(s,t)
  &:=\boldsymbol r_\Gamma(s)+t\nu_\Gamma(s),
  &|t|&<\delta_{\mathrm{col}},
  \label{eq:interface-normal-neighborhood}
\end{align}
在该邻域内单射，并且任一点至多被 $N_{\mathrm{ov}}$ 个此类邻域覆盖。再要求
\begin{align}
  \vartheta_{\mathrm{col}}
  &:=1-\delta_{\mathrm{col}}\kappa_*>0,
  &c_0&:=\frac{N_{\mathrm{ov}}}
    {\vartheta_{\mathrm{col}}\delta_{\mathrm{col}}},
  &c_1&:=\frac{N_{\mathrm{ov}}\delta_{\mathrm{col}}}
    {2\vartheta_{\mathrm{col}}}.
  \label{eq:trace-geometry-constants}
\end{align}
这里 $\delta_{\mathrm{col}}$ 是几何法向邻域的半宽，不是第~\ref{sec:numerics}~节的
Nystr\"om 网格步长 $h$；$\vartheta_{\mathrm{col}}$ 是法向坐标 Jacobian 的下界，也不同于
式~\eqref{eq:background-bloch-fiber} 的 $x$ 向 Bloch 参数 $\theta$。单连通性本身并不保证
这些条件，计算可靠上界时还必须验证邻域单射性、不同界面及周期副本的分离性和覆盖重数界。

\begin{proposition}[全界面正移位迹估计]
  \label{prop:positive-shift-trace-bound}
  在上述几何条件以及 $0<\rho_-\leq\rho\leq\rho_+$ 下，对每个
  $v\in D(A(\beta))$，有
  \begin{align}
    T_{\ell_*}(v)^2
    &:={}
    \sum_{\Gamma\in\mathscr E}
    \left(
      \lVert v|_\Gamma\rVert_{L^2(\Gamma)}^2
      +\ell_*^2
       \lVert\partial_{\nu_\Gamma}v|_\Gamma\rVert_{L^2(\Gamma)}^2
    \right)
    \notag\\
    &\leq C_{\mathrm{tr},\ell_*}^2
      \lVert(A(\beta)+\mu I)v\rVert_{\mathcal H}^2,
    \label{eq:positive-shift-trace-bound}\\
    C_{\mathrm{tr},\ell_*}^2
    &:={}
    \frac{c_0}{\rho_-\mu^2}
    +\frac{c_1+\ell_*^2c_0}{4\mu}
    +\ell_*^2c_1\rho_+.
    \label{eq:trace-constant}
  \end{align}
\end{proposition}

\begin{proof}
  先说明所需的全局二阶控制。把 $v\in D(A(\beta))$ 以相位
  $\ee^{\ii\beta d}$ 沿 $y$ 延拓，并在 $y$ 方向用
  $\psi_m$、在 $x$ 方向用 Fourier 变换。由于 $v\in H^1(B)$ 且其分布 Laplacian 属于
  $L^2(B)$，Parseval 等式给出
  \begin{align*}
    \sum_m\int_{\RR}
    (\xi^2+\beta_m^2)^2|\widehat v_m(\xi)|^2\,\dd\xi<\infty.
  \end{align*}
  因而 $v\in H^2(B)$，并且
  $\lVert D^2v\rVert_{L^2(B)}=\lVert\Delta v\rVert_{L^2(B)}$。这一结论使用全局
  分布 Laplacian，故材料界面上的传输条件已经包含在 $D(A(\beta))$ 中。

  令 $f=(A(\beta)+\mu I)v$。由 $A(\beta)$ 的非负自伴性和谱定理，
  \begin{align*}
    \lVert v\rVert_{L^2(B)}^2
    &\leq\frac{1}{\rho_-\mu^2}\lVert f\rVert_{\mathcal H}^2,
    &\lVert\nabla v\rVert_{L^2(B)}^2
    &\leq\frac{1}{4\mu}\lVert f\rVert_{\mathcal H}^2,\\
    \lVert D^2v\rVert_{L^2(B)}^2
    &\leq\rho_+\lVert f\rVert_{\mathcal H}^2.
  \end{align*}
  最后一式来自 $\Delta v=-\rho A(\beta)v$、
  $\lVert A(\beta)v\rVert_{\mathcal H}\leq\lVert f\rVert_{\mathcal H}$ 及上述 Hessian
  恒等式。

  对 $g\in H^1(-\delta_{\mathrm{col}},\delta_{\mathrm{col}})$，一维基本定理和
  Cauchy--Schwarz 不等式给出
  \begin{align*}
    |g(0)|^2
    \leq\frac{1}{\delta_{\mathrm{col}}}
      \int_{-\delta_{\mathrm{col}}}^{\delta_{\mathrm{col}}}|g(t)|^2\,\dd t
    +\frac{\delta_{\mathrm{col}}}{2}
      \int_{-\delta_{\mathrm{col}}}^{\delta_{\mathrm{col}}}|g'(t)|^2\,\dd t.
  \end{align*}
  沿式~\eqref{eq:interface-normal-neighborhood} 的法线使用此式，利用 Jacobian 下界
  $\vartheta_{\mathrm{col}}$，再对任意有限界面子族求和并使用覆盖重数，得到
  \begin{align*}
    \sum_\Gamma\lVert v|_\Gamma\rVert_2^2
    &\leq c_0\lVert v\rVert_2^2+c_1\lVert\nabla v\rVert_2^2,\\
    \sum_\Gamma\lVert\partial_{\nu_\Gamma}v|_\Gamma\rVert_2^2
    &\leq c_0\lVert\nabla v\rVert_2^2+c_1\lVert D^2v\rVert_2^2.
  \end{align*}
  对递增有限子族取极限并代入前三个体范数估计，恰得
  式~\eqref{eq:positive-shift-trace-bound}--\eqref{eq:trace-constant}。
\end{proof}

后文简记
\begin{align}
  C&:=C_{\mathrm{tr},\ell_*},
  &Q&:=C\varepsilon_{\partial,\ell_*}.
  \label{eq:trace-defect-product}
\end{align}
$C$ 只依赖式~\eqref{eq:trace-constant} 中同一个 $\mu$、材料上下界、参考长度和经过验证的
几何输入；它不是从候选矩阵残量拟合出来的量。

\subsection{谱距离定理与可计算的可靠上界}

\begin{theorem}[原始界面缺陷的谱距离界]
  \label{thm:resolvent-spectral-distance}
  设 $A(\beta)$ 是式~\eqref{eq:fixed-beta-operator} 的非负自伴算子，界面族满足
  命题~\ref{prop:positive-shift-trace-bound} 的几何条件。设非零原始场
  $u_{h,M}\in\mathcal H$ 逐片属于 $H^2$，满足
  式~\eqref{eq:raw-piecewise-pde}、精确的 $\beta$-准周期条件，并且
  式~\eqref{eq:boundary-defect-norm} 有限。若
  \begin{align}
    Q&<U,
    \label{eq:positive-resolvent-denominator}
  \end{align}
  则
  \begin{align}
    \operatorname{dist}\!\left(\mu,\sigma(A(\beta))\right)
    &\leq\frac{2\mu Q}{U-Q}.
    \label{eq:resolvent-spectral-distance-bound}
  \end{align}
  这个结论是到整个算子谱的距离上界，不要求预先知道本质谱隙。
\end{theorem}

\begin{proof}
  若 $\mu\in\sigma(A(\beta))$，左端为零，结论成立。以下设
  $\mu\notin\sigma(A(\beta))$，并只为证明定义辅助场
  \begin{align}
    z&:=(A(\beta)-\mu I)^{-1}u_{h,M}
    \in D(A(\beta)).
    \label{eq:auxiliary-resolvent-field}
  \end{align}
  谱距离估计的计算不需求解 $z$。

  先证明所需的 Green 恒等式。在每个材料分片上对
  $u_{h,M}$ 和 $z$ 使用第二 Green 恒等式。共享界面两侧的外法向分别为
  $\nu_\Gamma$ 和 $-\nu_\Gamma$，而 $z\in D(A(\beta))\subset H^2(B)$ 的值迹和普通
  法向导数迹均单值。结合式~\eqref{eq:raw-interface-jumps}，两侧边界项之和恰为
  \begin{align*}
    j_\Gamma\overline z
    +a_\Gamma\partial_{\nu_\Gamma}\overline z.
  \end{align*}
  为处理无穷远，取实值截断函数 $\chi_R(x)$，使其在 $|x|\leq R$ 时为一、在
  $|x|\geq2R$ 时为零，并满足
  $\lVert\chi_R'\rVert_\infty=O(R^{-1})$、
  $\lVert\chi_R''\rVert_\infty=O(R^{-2})$。以 $\chi_Rz$ 作第二个场并先对有限个分片求和。
  顶、底边界项由精确准周期条件相消；含 $\chi_R'$、$\chi_R''$ 的体项由
  $u_{h,M}\in\mathcal H$ 和 $z\in H^2(B)$ 控制并趋于零，界面上的附加截断项则由
  $\sum_\Gamma\lVert a_\Gamma\rVert_2^2<\infty$ 与命题
  ~\ref{prop:positive-shift-trace-bound} 控制并趋于零。主界面级数由
  式~\eqref{eq:boundary-defect-norm}、该命题和 Cauchy--Schwarz 不等式绝对收敛。因此令
  $R\to\infty$ 得到
  \begin{align}
    U^2
    =\sum_{\Gamma\in\mathscr E}
    \int_\Gamma
    \left(
      j_\Gamma\overline z
      +a_\Gamma\partial_{\nu_\Gamma}\overline z
    \right)\,\dd s.
    \label{eq:raw-green-identity}
  \end{align}

  记
  $\delta=\operatorname{dist}(\mu,\sigma(A(\beta)))>0$。由自伴算子的谱定理，
  $\lVert(A(\beta)-\mu I)^{-1}\rVert=\delta^{-1}$。又由
  $(A(\beta)+\mu I)z=u_{h,M}+2\mu z$，
  \begin{align}
    \lVert(A(\beta)+\mu I)z\rVert_{\mathcal H}
    &\leq\left(1+\frac{2\mu}{\delta}\right)U.
    \label{eq:positive-shift-of-resolvent}
  \end{align}
  对式~\eqref{eq:raw-green-identity} 使用带 $\ell_*$ 权重的 Cauchy--Schwarz 不等式和
  命题~\ref{prop:positive-shift-trace-bound}，得到
  \begin{align*}
    U^2
    &\leq\varepsilon_{\partial,\ell_*}T_{\ell_*}(z)
    \leq Q\left(1+\frac{2\mu}{\delta}\right)U.
  \end{align*}
  除以 $U>0$，再利用 $U-Q>0$ 移项，即有
  $\delta(U-Q)\leq2\mu Q$，这就是式
  ~\eqref{eq:resolvent-spectral-distance-bound}。
\end{proof}

用数值近似评价上述理论公式可得到谱距离估计；若要得到可靠上界，则必须为定理输入建立
可靠单侧界。具体地，设
\begin{align}
  C&\leq\overline C,
  &\varepsilon_{\partial,\ell_*}&\leq\overline\varepsilon_\partial,
  &0<\underline U&\leq U,
  \label{eq:certified-one-sided-inputs}\\
  \overline Q&:=\overline C\,\overline\varepsilon_\partial,
  &\underline D&:=\underline U-\overline Q>0.
  \label{eq:certified-positive-denominator}
\end{align}
$\overline C$ 来自材料、候选 $\mu$ 和法向邻域几何的向上包围；
$\overline\varepsilon_\partial$ 来自同一原始场全部实际值跳跃和法向跳跃的向上包围；
$\underline U$ 则是该场加权质量的向下包围。三者不得混用不同候选、不同归一化或重新求解
后的密度。由定理~\ref{thm:resolvent-spectral-distance} 立即得到可计算的可靠上界
\begin{align}
  R_{\mathrm{spec}}
  &:=\frac{2\mu\overline C\,\overline\varepsilon_\partial}
    {\underline U-\overline C\,\overline\varepsilon_\partial},
  &\operatorname{dist}\!\left(\mu,\sigma(A(\beta))\right)
  &\leq R_{\mathrm{spec}}.
  \label{eq:certified-spectral-distance}
\end{align}
正分母条件~\eqref{eq:certified-positive-denominator} 是有限界成立的必要检查；若它不成立，
公式不提供有限谱距离界。$R_{\mathrm{spec}}$ 是到算子谱的距离上界，不是残量范数。
因此向外舍入的区间
\begin{align}
  I_{\mathrm{spec}}
  &:=[\mu-R_{\mathrm{spec}},\mu+R_{\mathrm{spec}}]
  \label{eq:certified-spectral-interval}
\end{align}
必与 $\sigma(A(\beta))$ 相交。这个结论既不识别最近谱点的编号或重数，也不保证区间内只有
一个谱点。

若另行证明某个开区间 $G$ 与 $\sigma_{\mathrm{ess}}(A(\beta))$ 不交，并且
$I_{\mathrm{spec}}\subset G$，则
命题~\ref{prop:discrete-spectrum-in-gap} 才允许把其中至少一个谱点解释为孤立、有限重的
导模特征值。本文的谱距离可靠上界本身不证明这样的 $G$ 存在，也不把猜想
~\ref{conj:essential-spectrum} 当成数值带隙的证明。

\paragraph{体缺陷的独立扩展。}
当前连续核原始场逐片精确满足式~\eqref{eq:raw-piecewise-pde}，故体缺陷为零。若将来使用
不精确满足 PDE 的分片场，令
$r=(-\Delta-\mu\rho)u_{h,M}$ 且
$\varepsilon_{\mathrm{vol}}:=\lVert\rho^{-1}r\rVert_{\mathcal H}$，则同一证明在
式~\eqref{eq:raw-green-identity} 中增加体积分，并给出
\begin{align*}
  \operatorname{dist}\!\left(\mu,\sigma(A(\beta))\right)
  \leq
  \frac{2\mu\overline C\,\overline\varepsilon_\partial
    +\overline\varepsilon_{\mathrm{vol}}}
  {\underline U-\overline C\,\overline\varepsilon_\partial}.
\end{align*}
这一扩展不参与本文算例，也不改变精确 PDE 主线。


\clearpage
\appendix

\section{中心胞元表示定理的证明}
\label{app:center-cell-representation}

\begin{proof}[定理~\ref{thm:center-cell-representation} 的证明]
  在每个 $\partial\Omega_{0,\ell}$ 上，以下用上标 $+$ 表示从介质柱外部取得的迹，
  用上标 $-$ 表示从介质柱内部取得的迹；这两个上标与左右人工边界
  $\Gamma_-$、$\Gamma_+$ 无关。法向 $\nu$ 始终由介质柱内部指向外部。

  先确定均匀背景场。由于介质柱与 $\Gamma_\pm$ 严格分离，$u$ 在两个人工边界的
  邻域内满足波数为 $\kappa_{\mathrm e}$ 的均匀 Helmholtz 方程。将 $u$ 及
  $\partial_xu$ 的迹投影到 $\psi_m$ 后，每个 Fourier 分量的两个 Rayleigh 系数由
  \begin{equation*}
    \begin{bmatrix}
      1&1\\ \ii\gamma_m&-\ii\gamma_m
    \end{bmatrix}
  \end{equation*}
  所给出的二元关系确定。因为 $\gamma_m\neq0$，该矩阵可逆，所以锚定在
  $\Gamma_-$ 的右行系数 $a_{0,m}^-$ 与锚定在 $\Gamma_+$ 的左行系数
  $b_{0,m}^+$ 均由 $u$ 唯一确定。用这些系数构造正文中的
  $u_{\mathrm{hom}}$，并在中心胞元的背景区域令
  \begin{equation*}
    w^+:=u-u_{\mathrm{hom}}.
  \end{equation*}
  按构造，$w^+$ 在左端没有右行入射分量，在右端没有左行入射分量，因而在两个端口都满足
  outgoing Rayleigh 条件。对第 $\ell$ 根介质柱，再令 $w_\ell^-:=u|_{\Omega_{0,\ell}}$。

  接着构造层势密度。对每个 $\ell$，考虑交换内外波数的辅助问题：在
  $\Omega_{0,\ell}$ 内使用 $\kappa_{\mathrm e}$，在
  $\RR^2\setminus\overline{\Omega_{0,\ell}}$ 内使用
  $\kappa_{\mathrm i,\ell}$，并要求外部场满足 outgoing Sommerfeld 条件。其界面跳跃取为
  \begin{align*}
    v_\ell^+-v_\ell^-
    &=-w^+-w_\ell^-,\\
    \partial_\nu v_\ell^+-\partial_\nu v_\ell^-
    &=-\partial_\nu w^+-\partial_\nu w_\ell^-.
  \end{align*}
  定理的非例外性假设保证该问题具有唯一解 $v_\ell$。定义
  \begin{align*}
    \tau_\ell
    &:=w^+-v_\ell^-
      =-w_\ell^--v_\ell^+,\\
    \sigma_\ell
    &:=-\partial_\nu w^++\partial_\nu v_\ell^-
      =\partial_\nu w_\ell^-+\partial_\nu v_\ell^+.
  \end{align*}
  两组等式分别由辅助问题的两个跳跃条件得到。

  现在验证这些密度给出所需场。自由空间 Green 表示公式及其消去公式表明，对
  $\mathbf z\in\Omega_{0,\ell}$，
  \begin{equation*}
    \left(
      \mathcal S_\ell^{(\kappa_{\mathrm i,\ell})}[\sigma_\ell]
      +\mathcal D_\ell^{(\kappa_{\mathrm i,\ell})}[\tau_\ell]
    \right)(\mathbf z)
    =w_\ell^-(\mathbf z).
  \end{equation*}
  事实上，代入 $\tau_\ell$、$\sigma_\ell$ 的第二组表达式后，来自
  $w_\ell^-$ 的两项由内部 Green 公式重构 $w_\ell^-$，来自外部场
  $v_\ell^+$ 的两项则在介质柱内部相消。

  对背景区域，在
  $\mathcal C_0\setminus\overline{\Omega_0}$ 上对 $w^+$ 与
  $G_{\mathrm{QP}}^{(\kappa_{\mathrm e})}$ 使用 Green 第二恒等式。顶、底边界积分因
  $\beta$-准周期相位互相抵消；两个竖直端口上的积分在逐项代入 outgoing Rayleigh 展开后
  也为零。这里每个 Fourier 分量的消去因子是同一个 $\gamma_m$，而
  $\gamma_m\neq0$ 保证系数分解没有退化。因此
  \begin{equation*}
    \sum_{\ell=1}^{N_0}\left(
      -\mathcal S_{\mathrm{QP},\ell}^{(\kappa_{\mathrm e})}
        [\partial_\nu w^+]
      +\mathcal D_{\mathrm{QP},\ell}^{(\kappa_{\mathrm e})}[w^+]
    \right)
    =w^+.
  \end{equation*}
  对 $v_\ell^-$ 使用内部 Green 公式及其在介质柱外部的消去公式，则有
  \begin{equation*}
    \mathcal S_{\mathrm{QP},\ell}^{(\kappa_{\mathrm e})}
      [\partial_\nu v_\ell^-]
    -\mathcal D_{\mathrm{QP},\ell}^{(\kappa_{\mathrm e})}[v_\ell^-]
    =0
    \qquad\text{在 }\mathcal C_0\setminus\overline{\Omega_0}.
  \end{equation*}
  将 $\tau_\ell$、$\sigma_\ell$ 的第一组表达式代入并对 $\ell$ 求和，便得到
  \begin{equation*}
    \sum_{\ell=1}^{N_0}\left(
      \mathcal S_{\mathrm{QP},\ell}^{(\kappa_{\mathrm e})}[\sigma_\ell]
      +\mathcal D_{\mathrm{QP},\ell}^{(\kappa_{\mathrm e})}[\tau_\ell]
    \right)
    =w^+.
  \end{equation*}
  再加回 $u_{\mathrm{hom}}$，即得定理所述的外部表示；结合前面的内部表示，存在性得证。

  最后讨论唯一性。两组表示之差仍具有定理中的形式，并表示零场。前述端口系数提取关系立即
  给出所有 $a_{0,m}^-$、$b_{0,m}^+$ 均为零，所以 Rayleigh 系数列无条件唯一。若再满足定理
  末尾的互补系统零核条件，则只剩的层势部分在交换适用区域并改变内部层势符号后，由层势
  跳跃关系给出该齐次互补系统的一个解；它在每个界面上的共同 Dirichlet 和 Neumann 迹分别
  为 $-\tau_\ell$ 和 $\sigma_\ell$。零核条件迫使互补场恒为零，再次使用层势跳跃关系便有
  $\tau_\ell=0$、$\sigma_\ell=0$，$1\leq\ell\leq N_0$。因此在这一附加条件下，全部层势
  密度也唯一。
\end{proof}


\clearpage
\bibliographystyle{unsrt}
\bibliography{../reference}

\end{document}
